\centering
    \begin{tabular}{|c|c|c|}
      \hline
      Ток накала, А & $k$, мА/В$^{3/2}$ & $\sigma_k$, мА/В$^{3/2}$ \\
      \hline
      1.3 & $1.455 \times 10^{-2}$ & $1.0 \times 10^{-4}$ \\
      1.4 & $1.529 \times 10^{-2}$ & $0.9 \times 10^{-4}$ \\
      1.5 & $1.564 \times 10^{-2}$ & $1.1 \times 10^{-4}$ \\
      1.6 & $1.594 \times 10^{-2}$ & $1.2 \times 10^{-4}$ \\
      \hline
    \end{tabular}
    \caption{Коэффициенты наклона}
    \label{tab:k}
  \end{table}
  
  \subsection{Расчёт удельного заряда электрона}
  
  По формуле
  \[
  \frac{e}{m} = \frac{1}{2} \left( \frac{k \cdot 10^{-3}}{\frac{4}{9} \varepsilon_0 \frac{2\pi l}{r_a} \alpha} \right)^2
  \]
  (здесь $k$ переведён в А/В$^{3/2}$) были вычислены значения $e/m$ для каждого тока накала:
  
  \begin{table}[h]
    \centering
    \begin{tabular}{|c|c|c|}
      \hline
      Ток накала, А & $e/m$, Кл/кг & $\Delta (e/m)$, Кл/кг \\
      \hline
      1.3 & $1.72 \times 10^{11}$ & $0.02 \times 10^{11}$ \\
      1.4 & $1.90 \times 10^{11}$ & $0.02 \times 10^{11}$ \\
      1.5 & $1.99 \times 10^{11}$ & $0.03 \times 10^{11}$ \\
      1.6 & $2.07 \times 10^{11}$ & $0.03 \times 10^{11}$ \\
      \hline
    \end{tabular}
    \caption{Результаты расчёта $e/m$}
    \label{tab:em}
  \end{table}
  
  Среднее значение:
  \[
  \langle e/m \rangle = (1.92 \pm 0.16) \times 10^{11} \text{ Кл/кг}.
  \]
  
  Табличное значение $e/m = 1.758820 \times 10^{11}$ Кл/кг. Полученный результат находится в пределах погрешности эксперимента.
  
  \section{Выводы}
  
  В ходе работы экспериментально подтверждён закон «трёх вторых» для цилиндрического вакуумного диода в области пространственного заряда при достаточно больших напряжениях ($U \ge 5$ В). При малых напряжениях наблюдается отклонение, вызванное начальной тепловой скоростью электронов. Определён удельный заряд электрона:
  \[
  e/m = (1.92 \pm 0.16) \times 10^{11} \text{ Кл/кг},
  \]
  что хорошо согласуется с табличным значением $1.758820 \times 10^{11}$ Кл/кг. Отклонение может быть связано с неточностью определения коэффициента $\alpha$, погрешностью измерений, а также с неучётом начальных скоростей электронов при малых напряжениях.
  
\end{document}